\section{Related Work}
\label{sec:related_work}

Methods for preventing supervised learning algorithms from learning false associations between target variables and nuisance factors have been studied from various perspectives including ``feature selection''~\cite{bib:feat_select}, ``robustness through data augmentation''~\cite{bib:data_aug_1,bib:data_aug_2} and ``invariance induction''~\cite{bib:drl,bib:vfae,bib:sai}. Feature selection has typically been employed when data is available as a set of conceptual features, some of which are irrelevant to the prediction tasks. Our approach can be interpreted as an implicit feature selection mechanism for neural networks, which can work on both raw data (such as images) and feature-sets (e.g., frequency features computed from raw text). Popular feature selection methods~\cite{bib:feat_select} incorporate information-theoretic measures or use supervised methods to score features with their importance for the prediction task and prune the low-scoring ones. Our framework performs this task implicitly on latent features that the model learns by itself from the provided data.

Deep neural networks (DNNs) have outperformed traditional methods at several supervised learning tasks. However, they have a large number of parameters that need to be estimated from data, which makes them especially vulnerable to learning relationships between target variables and nuisance factors and, thus, overfitting. The most popular approach to expand the data size and prevent overfitting in deep learning has been synthetic data augmentation~\cite{bib:drl,bib:face_rec,bib:sbmrinet,bib:data_aug_1,bib:data_aug_2}, where multiple copies of data samples are created by altering variations of certain known nuisance factors. DNNs trained with data augmentation have been shown to generalize better and be more robust compared to those trained without in many domains including vision, speech and natural language. This approach works on the principle of inclusion. More specifically, the model learns to associate multiple seen variations of those nuisance factors to each target value. In contrast, our method encourages exclusion of information about nuisance factors from latent features used for predicting the target, thus creating more robust features. Furthermore, combining our method with data augmentation further helps our framework remove information about nuisance factors used to synthesize additional data, without the need to explicitly quantify or annotate the generated variations. This is especially helpful in cases where augmentation is performed using sophisticated analytical or composite techniques.

Several supervised methods for invariance induction and invariant feature learning have been developed recently, such as Controllable Adversarial Invariance (CAI)~\cite{bib:sai}, Variational Fair Autoencoder (VFAE)~\cite{bib:vfae}, and a maximum mean discrepancy based model (NN+MMD)~\cite{bib:nnmmd}. These methods use annotated information about variations of specific nuisance factors to force their exclusion from the learned latent representation. They have also been applied to learn ``fair'' representations based on domain knowledge, such as making predictions about the savings of a person invariant to age, where making the prediction task invariant to such factors is of higher priority than the prediction performance itself~\cite{bib:sai}. Our method induces invariance to nuisance factors with respect to a supervised task in an unsupervised way. However, it is not guaranteed to work in ``fairness'' settings because it does not incorporate any external knowledge about factors to induce invariance to.

Disentangled representation learning is closely related to our work since disentanglement is one of the pillars of invariance induction in our framework as the model learns two embeddings (for any given data sample) that are expected to be uncorrelated to each other. Our method shares some properties with multi-task learning (MTL)~\cite{bib:mtl} in the sense that the model is trained with multiple objectives. However, a fundamental difference between our framework and MTL is that the latter promotes a shared representation across tasks whereas the only information shared \emph{loosely} between the tasks of predicting $y$ and reconstructing $x$ in our framework is a noisy version of $e_1$ to help reconstruct $x$ when combined with a separate encoding $e_2$, where $e_1$ itself is used directly to predict $y$.